\begin{tikzpicture}[font=\sffamily\large,text=ink,
 labelnode/.style={align=center,inner sep=3pt},
 title/.style={font=\sffamily\bfseries\large,align=center,inner xsep=9pt,inner ysep=3pt},
 flow/.style={draw=ink!85,line width=.85pt,-{Stealth[length=1.55mm]},rounded corners=2pt},
 stage/.style={draw=sea,line width=1.25pt,-{Stealth[length=1.75mm]}},
 tag/.style={font=\sffamily\normalsize,inner sep=2pt}]
% The stage bands carry the sequence; illustrations belong to their own stage.
\filldraw[draw=sea!30,fill=sea!6] (0,0) rectangle (15,-2.08);
\filldraw[draw=sea!30,fill=sea!4] (0,-2.94) rectangle (15,-5.05);
\filldraw[draw=ochre!40,fill=ochre!6] (0,-5.91) rectangle (15,-14.30);
\node[title,fill=sea!18] at (7.5,0) {房间与火源参数};
\node[title,fill=sea!18] at (7.5,-2.94) {任务 1\quad 丙烷火验证};
\node[title,fill=ochre!21] at (7.5,-5.91) {任务 2\quad 丙烯火排烟设计与优化};
\draw[stage] (7.5,-2.12)--(7.5,-2.61);
\draw[stage] (7.5,-5.09)--(7.5,-5.58);
% Physical plan and door states, simplified from the report's layout.
\begin{scope}[shift={(.66,-1.18)},x=.52cm,y=.52cm]
\filldraw[draw=ink,fill=white,thick] (0,0) rectangle (6.6,.9);
\foreach \x in {.825,2.475,4.125,5.775}
 \draw[sea,line width=2pt] (\x-.215,.9)--(\x+.215,.9);
\fill[clay] (3.175,.325) rectangle (3.425,.575);
\end{scope}
\node[labelnode] at (2.38,-1.66) {房间长度 6.60 m};
\foreach \y/\id in {-.74/OP-1,-1.16/OP-6}{
 \node[anchor=east,font=\sffamily\normalsize] at (6.83,\y) {\id};
 \foreach \k in {0,...,3}{\draw[ink!65,fill=white] (7+.42*\k,\y-.1) rectangle ++(.28,.20);}
}
\foreach \k in {0,...,3}{\fill[sea] (7+.42*\k,-.84) rectangle ++(.28,.20);}
\fill[sea] (7+.42*2,-1.26) rectangle ++(.28,.20);
\node[labelnode] at (7.50,-1.66) {侧门开启状态};
% Two burner symbols indicate the two prescribed heat release rates.
\foreach \x/\s in {11.73/.68,12.82/1}{
 \begin{scope}[shift={(\x,-1.24)},scale=\s]
 \draw[ink,line width=1.2pt] (-.31,0)--(.31,0);
 \fill[clay] (0,.04).. controls (-.53,.17) and (-.13,.51) ..(-.10,.65)
 .. controls (.05,.53) and (.11,.42) ..(.14,.30)
 .. controls (.23,.45) and (.23,.47) ..(.24,.52)
 .. controls (.50,.18) and (.26,.06) ..(0,.04);
 \end{scope}}
\node[labelnode] at (12.32,-1.66) {47.2 / 106.2 kW};
% Centre thumbnail uses the actual paired profile for OP-1 at 106.2 kW.
\begin{scope}[shift={(1.57,-4.21)},x=.33cm,y=.21cm]
\draw[step=1,sea!65,thin] (0,0) grid (5,2);
\draw[ink] (0,0) rectangle (5,2);
\draw[ink] (0,2)--(.7,2.85)--(5.7,2.85)--(5,2) (5.7,2.85)--(5.7,.85)--(5,0);
\end{scope}
\node[labelnode] at (2.52,-4.64) {运行四组 FDS};
\begin{axis}[at={(6.43cm,-4.25cm)},anchor=south west,scale only axis,width=2.08cm,height=.77cm,
 axis lines=left,axis line style={-,ink!55},xtick=\empty,ytick=\empty,xmin=0,xmax=660,ymin=250,ymax=1000]
\addplot[sea,thick] table[col sep=comma,x=sensor_x_cm,y=fds_rise_K]{../../05_Calculations/op1_106_coarse_90_120.csv};
\addplot[clay,only marks,mark=*,mark size=.85pt] table[col sep=comma,x=sensor_x_cm,y=exp_rise_K]{../../05_Calculations/op1_106_coarse_90_120.csv};
\end{axis}
\node[labelnode] at (7.48,-4.64) {对照 51 个测点};
\node[labelnode] at (12.45,-4.12) {比较不同门态下的\\峰值与温度剖面};
\draw[flow] (4.52,-4.07)--(5.59,-4.07);
\draw[flow] (9.33,-4.07)--(10.39,-4.07);

% The left lane follows physical design decisions. The right lane analyses FDS records.
\tikzset{designnode/.style={labelnode,font=\sffamily\small},
 decision/.style={designnode,draw=ochre!80,fill=white,rounded corners=2pt,inner sep=4pt}}
\draw[ochre!32] (8.42,-6.35)--(8.42,-13.97);
\node[designnode] (initial) at (2.23,-6.98) {\textbf{16 个 OP-6 方案}\\达标数 0/16};
\node[designnode] (layout) at (6.28,-6.98) {\textbf{调整布局}\\垂壁、排烟与补风};
\draw[flow] (initial.east)--(layout.west);
\node[designnode] (coarse) at (6.28,-8.47) {\textbf{FDS 粗网格筛选}\\0.96、1.28 通过};
\draw[flow] (layout.south)--(coarse.north);
\node[designnode] (refine) at (2.23,-8.47) {\textbf{加密竖向网格}\\复查净高};
\draw[flow] (coarse.west)--(refine.east);
\node[decision] (gate) at (2.23,-10.05) {\textbf{测试 1.12}\\OP-6，烟炱产率 0.02};
\draw[flow] (refine.south)--(gate.north);
\node[designnode] (six) at (6.28,-10.05) {\textbf{复核 1.12}\\两种门态，三种产率};
\draw[flow] (gate.east)--(six.west);
\node[designnode] (verified) at (2.23,-11.85) {\textbf{八个测柱，120 s}\\$h_{\min}=0.4358$ m};
\draw[flow] (six.south)--(six.south |- verified.east)--(verified.east);
\node[designnode,fill=white,draw=ochre!55,inner sep=5pt] (power) at (4.16,-13.42)
 {$\displaystyle\min P$\quad s.t.\ 各测点净高 $\geq0.38$ m\\
 \textbf{选择 1.12 m$^3$/s}\quad 6.91 W；1.28 为 10.32 W};
\draw[flow] (verified.south)--(verified.south |- power.north);
\node[designnode,text=sea] (replay) at (11.82,-6.67) {\textbf{回放同组 16 个方案}};
\draw[flow] (initial.north)--++(0,.30)-|(8.72,-6.67)--(replay.west);
\foreach \x/\heading/\ranks in {
 10.25/{贝叶斯优化}/{1,9,12,2,10,16,5,11,14,6,15,7,3,13,8,4},
 13.22/{遗传算法}/{1,13,8,2,14,10,9,5,15,16,7,6,3,12,11,4}}{
 \node[align=center,font=\sffamily\normalsize] at (\x,-7.36) {\heading};
 \foreach \q [count=\i from 0] in \ranks{
  \pgfmathtruncatemacro{\rownum}{int(\i/4)}
  \pgfmathtruncatemacro{\colnum}{mod(\i,4)}
  \begin{scope}[shift={(\x+.42*\colnum-.63,-7.99-.34*\rownum)}]
   \ifnum\i=11
    \fill[sea] (-.21,-.17) rectangle (.21,.17);
    \node[text=white,font=\sffamily\bfseries\small] {\q};
   \else
    \filldraw[draw=ink!24,fill=white,line width=.3pt] (-.21,-.17) rectangle (.21,.17);
    \node[font=\sffamily\small] {\q};
   \fi
  \end{scope}}
}

\node[designnode] at (11.82,-9.83) {查到最大余量\\固定 16 \quad BO 7 \quad GA 6--11};
\node[designnode] (hist) at (11.82,-11.04) {\textbf{已归档 FDS 最低净高}};
\node[designnode] (pred) at (11.82,-12.18) {\textbf{PINN 与回归}\\预测候选方案最低净高};
\node[designnode,text=sea] (test) at (11.82,-13.53) {以留出 FDS 结果\\检验预筛精度};
\draw[flow] (hist.south)--(pred.north);
\draw[flow] (pred.south)--(test.north);
\draw[flow,draw=sea] (pred.west)--(8.80,-12.18)--node[pos=.48,above,sloped,font=\sffamily\small,text=sea]{建议检查顺序}(8.80,-8.47)--(coarse.east);
\end{tikzpicture}
